\section*{Hoofdstuk \ref{ch:b1:fluid-envelopes-structure} --- Structuur van de atmosfeer en de oceaan}
% values: calc/ch26.py

\begin{solution}{exo:b1:fluid-envelopes-structure:1}
(a)~$H = R^*T/(M_0 g_0) = 8314.32 \times 250/(28.9644 \times 9.80665) =
\qty{7.32}{km}$. (b)~$e^{10/7.32} = 3.92$: de druk daalt tot ongeveer een kwart.
(c)~$H$ is evenredig met de temperatuur: bij \qty{216.65}{K} is het
\qty{6.34}{km}, tegen \qty{8.43}{km} bij \qty{288.15}{K}. Koude lucht is bij een gegeven
druk dichter, dus daalt de druk sneller met de hoogte.
\end{solution}

\begin{solution}{exo:b1:fluid-envelopes-structure:2}
Met $\Gamma_d = \qty{9.76}{K/km}$: \qtyrange{0}{0.3}{km}, $11.0 > 9.76$,
labiel; \qtyrange{0.3}{1.2}{km}, $9.7$, neutraal binnen de precisie van de
gegevens; \qtyrange{1.2}{1.5}{km}, een inversie; \qtyrange{1.5}{8}{km}, stabiel;
boven \qty{8}{km}, isotherm, stabiel. De inversie is het meest stabiel, aangezien
$\Gamma_d - \Gamma = 13.8$ daar het grootst is.
\end{solution}

\begin{solution}{exo:b1:fluid-envelopes-structure:3}
(a)~$\sigma = 1026.95 - 1000 = \qty{26.95}{kg.m^{-3}}$. (b)~B, het lichtere, ligt
boven A. (c)~Alle dichtheden van zeewater beginnen met dezelfde cijfers, 10; $\sigma$ houdt
alleen de cijfers over die van het ene water naar het andere veranderen.
\end{solution}

\begin{solution}{exo:b1:fluid-envelopes-structure:4}
(a)~$S = 0.25 \times 35.40 + 0.75 \times 34.30 = 34.575$; $\theta = 0.25 \times
15.0 + 0.75 \times 4.0 = \qty{6.75}{\celsius}$. (b)~Op het lijnstuk PQ, een kwart van
de weg van Q naar P (\cref{prop:b1:fluid-envelopes-structure:mixing}).
(c)~De mengsels van P, Q en R vullen de driehoek PQR. Schrijven we
$(34.60, 8.0) = a\,\mathrm P + b\,\mathrm Q + c\,\mathrm R$ met $a + b + c = 1$,
dan volgt $a = 0.345$, $b = 0.788$, $c = -0.133$: een negatief aandeel van R, dus ligt het
punt buiten de driehoek en kan geen mengsel van de drie het voortbrengen.
\end{solution}

\begin{solution}{exo:b1:fluid-envelopes-structure:5}
(a)~$\theta = 253.15 \times 2^{2/7} = \qty{308.6}{K}$. (b)~Op \qty{1000}{hPa},
$\theta = T = \qty{298.15}{K}$. De \omterm{def:b1:fluid-envelopes-structure:potential-temperature}{potentiële temperatuur} neemt naar boven toe, van
298.2 tot \qty{308.6}{K}: de laag is stabiel voor droge lucht. (c)~Zij zou aankomen met
haar \omterm{def:b1:fluid-envelopes-structure:potential-temperature}{potentiële temperatuur}, \qty{308.6}{K}, dat is \qty{35.4}{\celsius}, tien
graden warmer dan de lucht die er al is.
\end{solution}

\begin{solution}{exo:b1:fluid-envelopes-structure:6}
(a)~$T = 288.15 - 6.5 \times 5 = \qty{255.65}{K}$ en $p = \num{101325} \times
(255.65/288.15)^{5.256} = \qty{54020}{Pa}$, \qty{540.2}{hPa}. (b)~$H =
\qty{8.43}{km}$ en $p = \num{101325}\,e^{-5/8.43} = \qty{56010}{Pa}$. De echte
(standaard-)druk is lager: de lucht daar is kouder dan \qty{288.15}{K}, dus dichter, en de druk daalt sneller met de hoogte. (c)~$p = \num{22632}\,
e^{-9/6.34} = \qty{5475}{Pa}$, de waarde van
\cref{tab:b1:fluid-envelopes-structure:layers}.
\end{solution}

\begin{solution}{exo:b1:fluid-envelopes-structure:7}
(a)~$\beta\,\Delta S = \alpha\,\Delta T$ geeft $\Delta S = 1.67/7.61 = 0.22$.
(b)~$\Delta\rho = 1026.95 \times (-1.67\times10^{-4} \times 2 + 7.61\times10^{-4}
\times 0.3) = \qty{-0.11}{kg.m^{-3}}$: het water wordt lichter.
(c)~$\Delta\rho = 1028.11 \times (-0.53\times10^{-4} \times 2 + 7.85\times10^{-4}
\times 0.3) = \qty{+0.13}{kg.m^{-3}}$: dezelfde veranderingen maken het nu dichter, omdat opwarming de dichtheid van water nabij het vriespunt nauwelijks verandert.
\end{solution}

\begin{solution}{exo:b1:fluid-envelopes-structure:8}
(a)~$N^2 = (9.81/260) \times (9.76 - 4)\times10^{-3} = \qty{2.17e-4}{s^{-2}}$,
$N = \qty{0.0147}{s^{-1}}$, periode $2\pi/N = \qty{426}{s}$, ongeveer zeven
minuten. (b)~$N^2 = (9.81/270) \times 14.76\times10^{-3} = \qty{5.36e-4}{s^{-2}}$,
$N = \qty{0.0232}{s^{-1}}$, periode \qty{271}{s}. (c)~$N$ gaat naar nul en de
periode groeit zonder grens: de herstellende kracht verdwijnt in een neutrale laag.
\end{solution}

\begin{solution}{exo:b1:fluid-envelopes-structure:9}
(a)~Zij is dichter, $\sigma_\theta = 27.20$ tegen 26.30: het AAIW is zoeter maar veel
kouder. (b)~$\Delta\sigma = 1026.95 \times [-1.67\times10^{-4} \times
(4.31 - 15.43) + 7.61\times10^{-4} \times (34.30 - 35.54)] = 1.907 - 0.969 =
0.94$, tegen $0.90$ exact. (c)~Temperatuur: de afkoeling voegt
\qty{1.91}{kg.m^{-3}} toe, het zoeter worden haalt \qty{0.97}{kg.m^{-3}} weg.
\end{solution}

\begin{solution}{exo:b1:fluid-envelopes-structure:10}
(a)~Volgens \cref{thm:b1:fluid-envelopes-structure:stability}, $(g/\theta)\,\mathrm
d\theta/\mathrm dz = (g/T)(\Gamma_d - \Gamma)$; $\mathrm d\theta/\mathrm dz = 0$
geeft $\Gamma = \Gamma_d$. (b)~$T = T_s - \Gamma_d z$, een laag van constante gradiënt
$L = \Gamma_d$: volgens \cref{prop:b1:fluid-envelopes-structure:layer-pressure}, $p =
p_s (T/T_s)^{gM/(R\Gamma_d)}$, en $gM/(R\Gamma_d) = gMc_p/(Rg) = c_pM/R$
($= 3.5$ voor droge lucht). (c)~$p$ wordt nul waar $T = 0$, bij $z = T_s/\Gamma_d =
300/9.76 = \qty{30.7}{km}$. Werkelijke lagen van uniforme \omterm{def:b1:fluid-envelopes-structure:potential-temperature}{potentiële temperatuur} zijn hooguit een
paar kilometer diep: zij worden bedekt door stabiele lucht, zoals de troposfeer door
de stratosfeer.
\end{solution}

\begin{solution}{exo:b1:fluid-envelopes-structure:11}
(a)~$S = (33.787 + 36.050)/2 = 34.919$, $\theta = \qty{8.0}{\celsius}$.
(b)~Met een lineaire toestandsvergelijking is $\sigma$ een affiene functie van $S$ en
$\theta$, dus is de $\sigma$ van het mengsel het gemiddelde van de twee, 27.00.
(c)~$\alpha$ is \qty{0.75e-4}{K^{-1}} bij \qty{2}{\celsius} en
\qty{2.07e-4}{K^{-1}} bij \qty{14}{\celsius}, terwijl $\beta$ weinig verandert
(\num{7.80e-4} en \num{7.53e-4}). De dichtheid daalt steeds sneller naarmate de
temperatuur stijgt: zij is een concave functie van de temperatuur. Langs de mengingslijn ligt het mengsel op de gemiddelde temperatuur, waar de dichtheid boven de koorde
ligt die de twee ouders verbindt, dus boven 27.00. (d)~De overmaat komt van de kromming, een term van de tweede orde: zij groeit met het kwadraat van het temperatuurverschil van de ouders.
\end{solution}

\begin{solution}{exo:b1:fluid-envelopes-structure:12}
(a)~Het drukverschil is $5604 - 5500 = \qty{104}{dbar}$: het pakketje koelt af met
$0.129 \times 0.104 = \qty{0.013}{\celsius}$, tot \qty{0.718}{\celsius}.
(b)~Het water eromheen op \qty{5400}{m} is op \qty{0.730}{\celsius}, warmer: het pakketje, met bijna
dezelfde zoutgehalte, is dichter dan zijn omgeving en zakt terug. De laag is stabiel, hoewel de gemeten temperatuur naar beneden toeneemt; haar \omterm{def:b1:fluid-envelopes-structure:potential-temperature}{potentiële temperatuur} daalt naar beneden. (c)~$\theta \approx
0.7313 - 0.086 \times 5.604 = \qty{0.249}{\celsius}$, binnen een duizendste van een
graad van de exacte waarde.
\end{solution}

\begin{solution}{pb:b1:fluid-envelopes-structure:1}
\begin{enumerate}
  \item De temperatuur is bijna uniform tot \qty{50}{m}, daalt snel tot
        \qty{300}{m}, langzamer tot \qty{1000}{m}, en is onder die diepte bijna uniform.
        Het zoutgehalte daalt van 35.80 aan het oppervlak tot een minimum van 34.32 op
        \qty{800}{m}, en stijgt dan tot 34.90 op \qty{3000}{m}.
  \item $22.40 - 0.5 = \qty{21.90}{\celsius}$ wordt bereikt tussen \qty{50}{m}
        (22.10) en \qty{75}{m} (19.20): $50 + 25 \times 0.20/2.90 = \qty{52}{m}$.
  \item $(22.10 - 13.00)/250 = \qty{0.0364}{\celsius/m}$, \qty{3.6}{\celsius}
        per \qty{100}{m}. De \omterm{def:b1:fluid-envelopes-structure:ocean-layers}{thermoklien} strekt zich uit van de basis van de \omterm{def:b1:fluid-envelopes-structure:ocean-layers}{gemengde laag}, nabij \qty{50}{m}, tot ongeveer \qty{800}{m}; zij is het scherpst tussen
        \qty{50}{m} en \qty{75}{m} (\qty{11.6}{\celsius} per \qty{100}{m}).
  \item $\sigma = 1026.95 \times [1 - 1.67\times10^{-4} \times 12.4 +
        7.61\times10^{-4} \times 0.80] - 1000 = 25.45$, \qty{0.74}{kg.m^{-3}} te
        hoog. Het oppervlaktewater ligt ver van de referentietoestand, en $\alpha$
        is veel groter bij \qty{22}{\celsius} dan bij \qty{10}{\celsius}: de lineaire vorm onderschat de daling van de dichtheid door de warmte.
  \item Het zoutgehalteminimum op \qty{800}{m}, 34.32 bij \qty{4.8}{\celsius}, is
        Antarctisch Tussenwater, gevormd onder de regenrijke westenwinden van de
        Zuidelijke Oceaan. Het maximum op \qty{3000}{m}, 34.90 bij
        \qty{2.5}{\celsius}, is Noord-Atlantisch Diepwater, gevormd in de
        Groenland- en Labradorzee.
  \item $(2.50 - 2.26)/3.042 = \qty{0.079}{\celsius}$ per \qty{1000}{dbar}.
  \item $\sigma_\theta$ neemt met de diepte toe: de kolom is stabiel. De \omterm{def:b1:fluid-envelopes-structure:ocean-layers}{pyklien}
        is waar $\sigma_\theta$ het snelst toeneemt, in de bovenste paar honderd meter, samen met de \omterm{def:b1:fluid-envelopes-structure:ocean-layers}{thermoklien}: \num{0.82} tussen
        \qty{200}{m} en \qty{800}{m}, slechts \num{0.70} over de \qty{2200}{m} eronder.
  \item $S = (35.45 + 34.32)/2 = 34.885$; $\theta = (14.90 + 4.80)/2 =
        \qty{9.85}{\celsius}$.
  \item C: 26.461; I: 27.310; mengsel: 26.886, precies het gemiddelde van de twee
        (26.886), aangezien de lineaire vorm affien is.
  \item Het gemiddelde van de ouders is $(26.342 + 27.159)/2 = 26.751$; het mengsel,
        bij 26.888, is \qty{0.14}{kg.m^{-3}} dichter. Dit is cabbeling.
  \item Met de lineaire vorm geeft de worp 26.61 op \qty{300}{m}, 26.88 op
        \qty{500}{m} en 27.31 op \qty{800}{m}: de dichtheid van het mengsel ligt nabij
        \qty{500}{m}. Het water daar, $(34.82, \qty{9.6}{\celsius})$, is zeer
        dicht bij het mengsel $(34.885, \qty{9.85}{\celsius})$: de worp tussen het
        Central Water en het AAIW is zelf bijna een mengingslijn tussen hen.
  \item $S = 0.25 \times 35.45 + 0.75 \times 34.32 = 34.60$, $\theta = 0.25
        \times 14.90 + 0.75 \times 4.80 = \qty{7.33}{\celsius}$. Geïnterpoleerd
        tussen \qty{500}{m} en \qty{800}{m} geeft de worp
        $(34.57, \qty{7.2}{\celsius})$ nabij \qty{650}{m}: bijna erop.
  \item In \si{K/km}: A--B 10.0; B--C 9.7; C--D $-5.5$; D--E 6.29; E--F 6.50;
        F--G 6.57; G--H 0.33; H--I $-0.50$; I--J $-0.83$.
  \item A--B labiel; B--C neutraal; C--D een inversie; D--E, E--F, F--G stabiel;
        G--H bijna isotherm, stabiel; H--I en I--J inversies. De tropopauze
        ligt bij G, \qty{12.5}{km} en \qty{-59.0}{\celsius}, waar de temperatuur ophoudt te dalen.
  \item De gemiddelde temperatuur is \qty{214.65}{K}: $H = 287 \times 214.65/9.81 =
        \qty{6.28}{km}$, en $147\,e^{-4/6.28} = \qty{77.7}{hPa}$, in overeenstemming
        met de 78 gemeten.
  \item De lucht bij A, \qty{22.0}{\celsius}, is koeler dan de zee: warmte stroomt van de zee naar de lucht, die van onderen wordt verwarmd. De laag A--B, met een gradiënt die groter is dan $\Gamma_d$, is hiervan het resultaat.
  \item Bij B: $22.0 - 9.76 \times 0.5 = \qty{17.1}{\celsius}$, tegen 17.0; bij
        C: \qty{7.4}{\celsius}, tegen 7.3; bij D: \qty{3.5}{\celsius}, tegen
        9.5. Het pakketje is tot C iets warmer dan de lucht, en veel kouder bij D.
  \item Aan de basis van de inversie, nabij \qty{1.5}{km}: bij het binnengaan van de inversie wordt het pakketje kouder, dus dichter, dan de lucht eromheen, wat zijn stijging stopt.
  \item Pakketten die door de labiele oppervlaktelaag worden vrijgemaakt, stijgen erdoor en mengen haar tot haar \omterm{def:b1:fluid-envelopes-structure:potential-temperature}{potentiële temperatuur} uniform is, zodat haar gradiënt $\Gamma_d$ wordt; de inversie erboven stopt de menging.
  \item In K: A 294.15; B 293.99; C 294.08; D 300.51; E 312.65; F 329.29; G
        346.28; H 369.50; I 446.99; J 600.91.
  \item $\theta$ daalt van A naar B (labiel), is bijna constant van B naar C
        (neutraal) en stijgt daarboven (stabiel), het snelst van C naar D en boven
        G: dezelfde classificatie als in vraag~14.
  \item E--F: $(329.29 - 312.65)/4 = \qty{4.2}{K/km}$; H--I: $(446.99 -
        369.50)/4 = \qty{19.4}{K/km}$, bijna vijf keer zoveel.
  \item Met het gemiddelde $\theta = \qty{408.2}{K}$: $N^2 = (9.81/408.2) \times
        0.01937 = \qty{4.66e-4}{s^{-2}}$, $N = \qty{0.0216}{s^{-1}}$, een periode
        van \qty{291}{s}, ongeveer vijf minuten.
  \item $N^2 = (9.81/1027) \times (26.35 - 24.79)/150 = \qty{9.9e-5}{s^{-2}}$,
        $N = \qty{0.0100}{s^{-1}}$, een periode van \qty{630}{s}. De onderste stratosfeer is de stabielere, met $N^2$ bijna vijf keer groter
        dan in de \omterm{def:b1:fluid-envelopes-structure:ocean-layers}{thermoklien}.
  \item Bij G, $\theta = 214.15 \times (1000/186)^{2/7} = \qty{346}{K}$: de \omterm{def:b1:fluid-envelopes-structure:potential-temperature}{potentiële temperatuur} aan de tropopauze is ongeveer \qty{346}{K}.
\end{enumerate}
\end{solution}
